% #############################################################################
% Glossary entries, in alphabetical order
% !TEX root = ../../Main.tex
% #############################################################################
\newglossaryentry{alpha}
{
    name=alpha,
    description={The part of a strategy's mean return that its exposure to the market does not explain, taken from the regression of the strategy's returns on the market's returns}
}
\newglossaryentry{ask}
{
    name=ask price,
    description={The price at which the market sells, paid to open a long position or to close a short one}
}
\newglossaryentry{backtest}
{
    name=backtest,
    description={A simulation of a trading strategy on historical prices}
}
\newglossaryentry{balance}
{
    name=balance,
    description={The account value from closed positions only, before the result of any open position}
}
\newglossaryentry{base}
{
    name=base currency,
    description={The first currency of a pair, the one being bought or sold; the pair's price is the amount of quote currency one unit of it costs}
}
\newglossaryentry{beta}
{
    name=beta,
    description={The sensitivity of a strategy's return to the market's return, the slope of the regression of one on the other}
}
\newglossaryentry{bid}
{
    name=bid price,
    description={The price at which the market buys, received when closing a long position or opening a short one}
}
\newglossaryentry{calmar}
{
    name=Calmar ratio,
    description={The annualised return divided by the maximum drawdown}
}
\newglossaryentry{commission}
{
    name=commission,
    description={The fee a broker charges for the volume traded, on opening and on closing a position}
}
\newglossaryentry{pair}
{
    name=currency pair,
    description={The price of one currency expressed in another, such as EUR/USD}
}
\newglossaryentry{drawdown}
{
    name=drawdown,
    description={The fall of equity from its previous peak, as a fraction of that peak}
}
\newglossaryentry{equity}
{
    name=equity,
    description={The balance plus the unrealised result of the open positions}
}
\newglossaryentry{hedging}
{
    name=hedging account,
    description={An account that can hold several positions on one instrument, in either direction, at the same time}
}
\newglossaryentry{leverage}
{
    name=leverage,
    description={The ratio between the size of a position and the account capital that backs it}
}
\newglossaryentry{long}
{
    name=long position,
    description={A position that gains when the price rises}
}
\newglossaryentry{lot}
{
    name=lot,
    description={A standard trade size of 100,000 units of the base currency}
}
\newglossaryentry{netting}
{
    name=netting account,
    description={An account that holds at most one position per instrument, so an opposite order reduces, closes or reverses it}
}
\newglossaryentry{pip}
{
    name=pip,
    description={The fourth decimal place of most currency pair prices, and the second for pairs quoted in Japanese yen}
}
\newglossaryentry{quote}
{
    name=quote currency,
    description={The second currency of a pair, in which the price is expressed}
}
\newglossaryentry{sharpe}
{
    name=Sharpe ratio,
    description={The mean excess return divided by the standard deviation of returns}
}
\newglossaryentry{short}
{
    name=short position,
    description={A position that gains when the price falls}
}
\newglossaryentry{sortino}
{
    name=Sortino ratio,
    description={The mean excess return divided by the standard deviation of the negative returns only}
}
\newglossaryentry{spread}
{
    name=spread,
    description={The difference between the ask and the bid prices, paid on every round trip}
}
\newglossaryentry{swap}
{
    name=swap,
    description={The interest charged or paid for holding a position past the daily rollover time}
}
\newglossaryentry{tick}
{
    name=tick,
    description={A single change in the quoted ask or bid price}
}
\newglossaryentry{walkforward}
{
    name=walk-forward analysis,
    description={An evaluation in which a model is trained on one period and scored on the following one, repeated as the periods move forward in time}
}
\glsaddall